\documentclass[11pt,a4paper]{article}
\usepackage[utf8]{inputenc}
\usepackage[T1]{fontenc}
\usepackage{amsmath,amssymb,graphicx,booktabs,hyperref,listings,enumitem,xcolor}
\usepackage[margin=2.5cm]{geometry}
\usepackage{fancyhdr}
\pagestyle{fancy}\fancyhf{}\rhead{Fast Sorting}\cfoot{\thepage}
\graphicspath{{figs/}}
\newcommand{\R}{\mathbb{R}}
\newcommand{\bigO}[1]{\mathcal{O}(#1)}
\DeclareMathOperator{\argmin}{arg\,min}
\title{Fast Sorting of Nearly Sorted Data}
\author{Grace Hopper\thanks{Corresponding author.}\\Navy Research Lab \and John von Neumann\\IAS}
\date{6 October 2026}
\begin{document}
\maketitle
\begin{abstract}
We present an adaptive algorithm that sorts nearly sorted arrays in $\bigO{n \log k}$ time,
where $k$ is the number of inversions. Experiments in Section~\ref{sec:exp} confirm it.
\end{abstract}
\section{Introduction}\label{sec:intro}
Sorting is fundamental~\cite{knuth1998}. Our contributions are:
\begin{itemize}[noitemsep]
  \item an adaptive merge strategy (Section~\ref{sec:method});
  \item a tight bound, Equation~\eqref{eq:bound};
  \item an implementation in Python.\footnote{Available at \url{https://example.org/sort}.}
\end{itemize}
\section{Method}\label{sec:method}
Let $x \in \R^n$ and define the disorder
\begin{equation}\label{eq:bound}
  D(x) = \sum_{i<j} \mathbf{1}[x_i > x_j] \le \binom{n}{2}.
\end{equation}
We choose the run length $r^\star = \argmin_{r \ge 1} \left( \frac{n}{r} + r \log r \right)$.
\begin{align}
  T(n) &= 2\,T(n/2) + \bigO{n} \label{eq:rec}\\
       &= \bigO{n \log n} \notag
\end{align}
\begin{lstlisting}[language=Python]
def adaptive_sort(xs):
    runs = find_runs(xs)
    return merge_all(runs)
\end{lstlisting}
\section{Experiments}\label{sec:exp}
Table~\ref{tab:times} and Figure~\ref{fig:speed} summarise the results.
\begin{table}[t]
\centering
\caption{Running time in milliseconds.}\label{tab:times}
\begin{tabular}{lrr}
\toprule
Input & Timsort & Ours \\
\midrule
Sorted & 1.2 & \textbf{0.9} \\
Reversed & 1.4 & \textbf{1.1} \\
Random & 9.8 & 10.1 \\
\bottomrule
\end{tabular}
\end{table}
\begin{figure}[t]
\centering
\includegraphics[width=0.6\linewidth]{plot}
\caption{Speed-up over the baseline.}\label{fig:speed}
\end{figure}
\section{Conclusion}
The recurrence~\eqref{eq:rec} explains the speed-up; see also \autoref{sec:intro}.
\end{document}
